\section{Stagewise interfaces for the fixed-level refinements}
\label{sec:polar}

The arbitrary-redundancy theorem uses simultaneous contraction, not the
stagewise statements below.  The retained R6--R7 appendix uses them to obtain
sharper small-field bounds and isolate the first exceptional modular
components; the companion R8--R10 calculations are not claims of this
submission.  This section records the interfaces needed here, with proofs, so
that the submission is self-contained.

For \(\lambda\in E^\vee\), divided-power contraction satisfies
\begin{equation}\label{eq:lift}
 g\in W_{\iota_\lambda f}\quad\Longleftrightarrow\quad
 \lambda g\in W_f.
\end{equation}
The lifted form is squarefree precisely when \(g\) is squarefree and avoids
the marked root \(\lambda\).

\begin{theorem}[Polar-flag construction]\label{thm:polar-construction}%
\coverage{conditional_deduction}%
\lean{PRSPolarInduction.iteratedProjectiveSequenceContraction_map,sequenceContraction_agrees_with_finite,PointedKernelLift.lift_splitSquarefreeKernelMember}%
Iterated contraction commutes with field extension and
\(\operatorname{GL}(E)\).  If
\(g\in W_{\iota_{\lambda_j}\cdots\iota_{\lambda_1}f}\), then
\(\lambda_1\cdots\lambda_jg\in W_f\).  If the projective markers
\(\lambda_1,\ldots,\lambda_j\) are pairwise distinct, the lift is squarefree
exactly when \(g\) is squarefree and avoids every marker.
\end{theorem}

\begin{proof}
Contraction is defined by the perfect divided-power/symmetric-power pairing,
so it commutes with base change and change of basis.  Repeated application of
\eqref{eq:lift} gives membership.  The product of the distinct markers is
squarefree, and its product with \(g\) is squarefree exactly when \(g\) is
squarefree and shares no marker factor.
\end{proof}

\begin{definition}[One-step lower package]\label{def:lower-package}
A one-step lower package consists of a closed lower bad locus, finitely many
geometric strata on its complement, a proper geometrically integral splitting
curve on each stratum with stated genus and rational-point bound, a
zero-dimensional deletion scheme containing branch, diagonal, collision, and
marker incidences, and a finite parameter scheme containing every unavailable
marker.  Every rational point outside the deletion scheme supplies a split
squarefree lower witness avoiding all retained markers.
\end{definition}

\begin{definition}[Contained-component assertion]\label{def:contained}
For a lower bad scheme \(B\) and a closed upper list \(\mathcal C\), the
contained-component assertion says that every syndrome whose contraction is
noninjective, whose entire polar line lies in \(B\), or whose collision divisor
is the whole marker line belongs to \(\mathcal C\).
\end{definition}

\begin{lemma}[Collision criterion]
\label{lem:marker-collision}%
\coverage{absent}%
If \(W_f\) has trivial gcd, then a marker \(\lambda\) is a repeated factor in
every lift through \(\lambda\) if and only if \(\lambda\) is a fixed factor of
\(W_{\iota_\lambda f}\); equivalently, it lies on the collision divisor of
the linear series \(W_f\).
\end{lemma}

\begin{proof}
Equation~\eqref{eq:lift} gives
\[
 \lambda W_{\iota_\lambda f}
 =W_f\cap\lambda\Sym^{n-2}E^\vee.
\]
After division by \(\lambda\), repeated lifting is therefore equivalent to a
fixed lower factor.  Since \(W_f\) is base-point-free, the same condition says
that the morphism defined by \(W_f\) has zero differential at \(\lambda\),
which is the collision condition.
\end{proof}

\begin{lemma}[Uniform collision bound]
\label{lem:uniform-collision}%
\coverage{absent}%
For \(j\geq5\), the collision divisor of a base-point-free
\(g^{\,j-4}_{\,j-2}\) on \(\PP^1\) has degree at most \(2j-6\) unless the
induced morphism is inseparable.  The only inseparable numerical cases are
\((j,p)=(5,3)\) and \((6,2)\).
\end{lemma}

\begin{proof}
An inseparable morphism factors through absolute Frobenius.  Hence
\(p\mid(j-2)\), while its \(j-3\) sections lie in the
\((j-2)/p+1\)-dimensional space of \(p\)-th powers; thus
\(p(j-4)\leq j-2\).  This leaves only \((j,p)=(5,3),(6,2)\).
In the separable case a nonzero value--derivative minor is a binary
Wronskian of degree \(2(j-2)-2=2j-6\) containing the collision scheme.
\end{proof}

\begin{theorem}[Finite-depth polar escape]
\label{thm:induction}%
\coverage{fragment}%
\lean{PRSPolarInduction.iteratedProjectiveSequenceContraction_map,sequenceContraction_agrees_with_finite,CoherentPolarInput.splitFree_implies_persistent_or_modular}%
Suppose \(s\) coherent contraction stages have unavailable-marker schemes of
degrees \(d_i<q+1\), each containing every previously chosen marker.  Suppose
each surviving terminal stratum has a proper
geometrically integral curve with
\[
 \#Y(\F_q)\geq q+\kappa-2g\sqrt q
\]
and at most \(\delta\) rational points that fail to supply a split squarefree
terminal witness avoiding all markers.  If
\[
 q+\kappa-2g\sqrt q>\delta,
\]
then \(W_f\) contains a split squarefree member.
\end{theorem}

\begin{proof}
Choose the markers successively.  A zero-dimensional scheme of degree below
\(q+1\) cannot contain every rational point of \(\PP^1\), so each stage has a
surviving choice.  On the terminal stratum the displayed inequality leaves a
rational point outside the bad set.  It gives a split squarefree
terminal member avoiding every marker.  Repeated use of \eqref{eq:lift}
multiplies by the pairwise distinct marker factors, and avoidance preserves
squarefreeness.
\end{proof}

\begin{proposition}[Contained rank-two polar lines]
\label{prop:contained-rank-two}%
\coverage{absent}%
Let \(n\geq5\), let \(f\in\PP(\Gamma^nE)\), and assume its contraction map is
injective.  Then the whole first-polar line lies in the lower rank-at-most-two
catalecticant locus if and only if \(f\) lies in the upper rank-at-most-two
locus.  Otherwise the rank-drop parameters form a scheme of degree at most
three on the marker line.
\end{proposition}

\begin{proof}
Let \(A=C_f^{(2)}\) and \(K=\ker A\).  Multiplication identifies
\(\Sym^{n-3}E^\vee\) with
\(H_\lambda=\lambda\Sym^{n-3}E^\vee\), and functoriality gives
\(C_{\iota_\lambda f}^{(2)}=A|_{H_\lambda}\).  The forward implication is
immediate.  Conversely, if \(\operatorname{rank}A=3\), then
\(\dim K=n-4\); lower rank at most two forces \(K\subset H_\lambda\).
Containment for every geometric \(\lambda\) would put \(K\) in the
intersection of all \(H_\lambda\), which is zero, a contradiction.
Finally, the maximal minors of \(A|_{H_\lambda}\) are cubics in \(\lambda\),
so in the noncontained case their common zero scheme has degree at most
three.
\end{proof}
